\chapter{Firmware interface and boot}\label{ch:firmware}

\section{The Strata Firmware Interface (SFI v1)}\label{sec:sfi}

The firmware is a \emph{library}, not a resident runtime. There is no firmware call table and no
firmware code running beneath the OS: the boot ROM includes the library (\texttt{firmware/fw.s}) in its
own image, and an OS that wants the same services includes it too, or replaces any routine with its
own. The contract consists of the API below and two platform facts: the machine describes itself in the
platform table, and devices are reached only through capabilities derived from the device root \reg{c3}.

\subsection{Conventions}
The firmware follows the calling convention of Appendix~\ref{app:abi}, with these additions:

\begin{table}[htbp]
\centering\small
\caption{SFI register, capability and tag conventions}\label{tab:sficonv}
\begin{tabularx}{\linewidth}{@{}l X@{}}
\toprule
\textbf{Resource} & \textbf{Use}\\\midrule
\reg{gp} & must hold the address of label \texttt{GPBASE}, which the including program defines at the
  start of its \texttt{.data}; firmware state is addressed relative to it\\
\reg{c0} & memory root, owned by the caller\\
\reg{c2} & console output (copyable)\\
\reg{c3} & device root (copyable)\\
\reg{c9} & buffer argument for block I/O\\
\reg{c10}, \reg{c11} & OS code and data shares, used by \texttt{fw\_start\_core}\\
\reg{c12}--\reg{c15} & clobbered by firmware routines; must not hold linear capabilities on entry
  (the routines write them, and would trap with \cause{capop})\\
save slots 62, 63 & block-device window, console window\\
\reg{tg12}--\reg{tg15} & used by firmware transfers\\
\reg{a0}--\reg{a7}, \reg{t0}--\reg{t6} & arguments, results and scratch; clobbered unless stated\\
\bottomrule
\end{tabularx}
\end{table}

\subsection{API}

\begin{table}[htbp]
\centering\small
\caption{SFI v1 routines}\label{tab:sfiapi}
\begin{tabularx}{\linewidth}{@{}>{\ttfamily}l X@{}}
\toprule
\textrm{\bfseries Routine} & \textbf{Description (result in \reg{a0})}\\\midrule
fw\_init() & read and verify the platform table (magic and checksum), save the console window; 0 or $-1$\\
fw\_cores() & number of cores, from the platform table\\
fw\_ndev() & number of platform table entries\\
fw\_dev(i) & D-LS address of entry $i$, or 0\\
fw\_find(type, n) & D-LS address of the $n$-th entry of a type (1 memory, 2 console, 3 platform table, 4 block device, 5 framebuffer, 6 timer), or 0\\
fw\_mem(space) & bytes of memory in a space (2 = T2, 3 = T3)\\
fw\_describe() & print the platform inventory, one line per entry\\
fw\_window(base, size) & \reg{c13} = a capability for device window $[\textit{base}, \textit{base}+\textit{size})$, cut from a copy of \reg{c3}\\
fw\_blk\_open(entry) & open a block device given its platform entry; 0 or $-1$\\
fw\_blk\_blocks() & capacity of the open block device, in blocks\\
fw\_blk\_read(lba, n, off) & read $n$ blocks into \reg{c9} at byte offset \textit{off}; status (Table~\ref{tab:blkstatus})\\
fw\_blk\_write(lba, n, off) & write $n$ blocks from \reg{c9} at byte offset \textit{off}; status\\
fw\_getc() & next console input byte; waits; 4 at end of input\\
fw\_puts(addr, len) & print \textit{len} bytes from D-LS address \textit{addr}\\
fw\_putc(ch), fw\_putnl() & print a character, a newline\\
fw\_putdec(v), fw\_puthex(v) & print unsigned decimal, or \texttt{0x} and 16 hex digits\\
fw\_start\_core(core) & give a secondary core the OS code share (\reg{c10}), data share (\reg{c11}) and the device root (\reg{c3}), in that order\\
\bottomrule
\end{tabularx}
\end{table}

\texttt{fw\_init}, \texttt{fw\_blk\_open}, the block routines and \texttt{fw\_getc} use the
capability save table, and the block routines take the buffer back with \insn{ctake}; they therefore run
only in supervisor mode, and a user-mode program that calls them traps with \cause{priv}. The printing
routines need only \reg{c2} and work in either mode. The firmware has no timer or framebuffer routines:
a program finds those devices with \texttt{fw\_find}, cuts a capability with \texttt{fw\_window} (or
with \insn{csplit} from its memory root, for the framebuffer) and uses plain transfers, as the
\texttt{clock} and \texttt{gfx} examples do (Chapter~\ref{ch:examples}).

The block routines follow the protocol of Section~\ref{sec:blkdev}: they build the command, submit it,
\insn{cgive} \reg{c9} to the device and take it back with \insn{ctake c9, t1}, accepting it only from
the device's node, then read the completion. Gifts other cores send meanwhile stay in their own queues. While the device
holds the buffer the caller cannot touch it, and when \insn{ctake} returns the data is in memory. The
printing routines return as soon as their \insn{tput} has issued; a later store into the same buffer
waits for it.

\section{Boot protocol}\label{sec:boot}

At reset every core runs the boot ROM image at I-LS address 0, in supervisor mode. The ROM's code and
data occupy the bottom of the Local Stores (under 2~KiB of each: it fits the smallest machine), and its
stack is where reset put \reg{sp}, the top of the D-LS. The ROM never leaves supervisor mode, and it
enters the OS in supervisor mode.

\begin{enumerate}
\item Core 0 prints a banner, calls \texttt{fw\_init} and \texttt{fw\_describe}, prints the Local Stores'
  sizes, finds block device 0 in the platform table and opens it.
\item It cuts its memory root into four pieces (Table~\ref{tab:memmap}): \reg{c0} keeps
  $[0, 15~\text{MiB})$; the load area $[15~\text{MiB}, 16~\text{MiB})$ goes to \reg{c9}; the memory above
  it, $[16~\text{MiB}, 48~\text{MiB})$, to \reg{c13}; and the framebuffer window
  $[48~\text{MiB}, 64~\text{MiB})$ to \reg{c14}. It reads block 0 into the load area and checks the boot
  header magic, and that the OS code (at I-LS \hex{4000}) and data (at D-LS \hex{A000}) fit the
  Local Stores the platform table describes.
\item It reads the OS code and data images from disk into the load area (at offsets \hex{10000} and
  \hex{20000}) and, if the header names one, the kernel's \term{far image} at offset \hex{30000} (at
  most \hex{D0000} bytes). It then \insn{csplit}s the load area into capabilities of exactly the images'
  sizes and \insn{cclear}s the rest.
\item It turns the images into read-only shares with \insn{cshare} (code in \reg{c10}, data in
  \reg{c11}, the far image, readable and executable, in \reg{c12}) and gives each secondary core a copy
  of the code and data shares, plus the device root, with \texttt{fw\_start\_core}. Secondary cores wait
  for these three gifts with \insn{ctake}. The far image share stays on core 0.
\item Every core copies the code into its own I-LS at \hex{4000} with \insn{tgetx} and the data into its
  D-LS at \hex{A000} with \insn{tget}, sets \reg{a0} = its core number, and jumps to \hex{4000} without
  waiting: fetching or loading bytes still in flight stalls.
\end{enumerate}

If a step fails (no valid platform table, no disk, a read error, no boot header, an OS that does not
fit), core 0 prints why, gives every other core a copy of its console capability \reg{c2} and halts.
A secondary core whose first gift is not executable (a code share is) halts too, so no core is left
waiting for the system core for good.

\begin{table}[htbp]
\centering\small
\caption{What the boot ROM hands to the OS}\label{tab:boothand}
\begin{tabularx}{\linewidth}{@{}l l X@{}}
\toprule
\textbf{Register} & \textbf{Cores} & \textbf{Capability}\\\midrule
\reg{c0} & core 0 & system memory $[0, 15~\text{MiB})$, \perm{RWX}\\
\reg{c1} & core 0 & the cluster store (untouched since reset)\\
\reg{c2} & every core & console output\\
\reg{c3} & every core & the device root\\
\reg{c4} & core 0 & core control (untouched since reset)\\
\reg{c10}, \reg{c11} & every core & read-only shares of the OS code and data images\\
\reg{c12} & core 0 & a read and execute share of the far image, exactly
  $[\hex{F30000}, \hex{F30000} + \textit{its length})$\\
\reg{c13} & core 0 & system memory $[16~\text{MiB}, 48~\text{MiB})$, \perm{RWX}\\
\reg{c14} & core 0 & the framebuffer window $[48~\text{MiB}, 64~\text{MiB})$, \perm{RWX}\\
\bottomrule
\end{tabularx}
\end{table}

The 1~MiB load area is never returned: it holds the OS images, now immutable. The far image holds code
the OS does not keep resident: its labels are far addresses relative to the start of \reg{c12}, and the
OS runs it with \insn{xcall} through \reg{c12} (Section~\ref{sec:os-kpaging}).

\subsection{Boot header}

\begin{table}[htbp]
\centering\small
\caption{Boot header (block 0, little-endian)}\label{tab:boothdr}
\begin{tabular}{@{}r l l@{}}
\toprule
\textbf{Offset} & \textbf{Type} & \textbf{Field}\\\midrule
0 & char[8] & magic \texttt{"STBOOT01"}\\
8 & u32 & code LBA\\
12 & u32 & code blocks\\
16 & u32 & code I-LS address (\hex{4000})\\
20 & u32 & data LBA\\
24 & u32 & data blocks\\
28 & u32 & data D-LS address (\hex{A000})\\
32 & u32 & entry I-LS address (\hex{4000})\\
44 & u32 & far image LBA (0 if there is none)\\
48 & u32 & far image length in blocks\\
56 & 4 bytes & \texttt{STDB}: the rest of the disk is a store (used by the OS)\\
60 & u32 & the store's page size in bytes (used by the OS)\\
\bottomrule
\end{tabular}
\end{table}

The boot ROM uses the LBA and block-count fields. The address fields document the layout; the
prototype ROM uses the fixed addresses \hex{4000} and \hex{A000} rather than reading them.

\subsection{Disk layout}

\texttt{mkdisk -o disk.img -k kernel.stx [-b blocks] [-p pagesize] [-B names] [[-t tags] file \ldots]}
(\texttt{tools/stdb/mkdisk.c}) builds a bootable disk image (default 32768 blocks, 16~MiB) from the
kernel image and any number of files. In pages of 4096 bytes (\texttt{-p}: 1024 or 2048):

\begin{center}\small
\begin{tabular}{@{}l l@{}}
\toprule
\textbf{Pages} & \textbf{Contents}\\\midrule
0 & boot header (its first block)\\
1, 2 & the store's two meta pages: a commit writes the one it did not write last\\
3\ldots & OS code image (I-LS bytes \hex{4000}--\hex{7FFF}), then OS data image (D-LS bytes \hex{A000}--\hex{DFFF})\\
then & the far image: the kernel image's T3 section (\texttt{.far}), whose offset 0 is far address 0 of \reg{c12}\\
rest & the store: copy-on-write B+trees (the file system's records, names, index and dictionary), the files' contents, free pages\\
\bottomrule
\end{tabular}
\end{center}

The code and data images are limited to 16~KiB each and the far image to \hex{D0000} bytes
(832~KiB). The files named by \texttt{-B} (by default \texttt{proc}, \texttt{dev} and \texttt{fs}) are
boot objects: each is stored in one run of pages and listed in the boot manifest, a page the newer
meta block points at, from which the kernel loads the first servers without reading a tree. The
format is described at the top of \texttt{tools/stdb/store.c}; \texttt{build/stdb check} verifies an
image.
